% mini_expose_v3_kurz.tex — Kurzfassung (max. 2 A4-Seiten)
% Kernaussage: Was will ich machen und warum.

\documentclass[11pt,a4paper]{scrartcl}
\usepackage[ngerman]{babel}
\usepackage[utf8]{inputenc}
\usepackage[T1]{fontenc}
\usepackage{lmodern}
\usepackage{csquotes}
\usepackage{microtype}
\usepackage{enumitem}
\usepackage{xcolor}
\usepackage{hyperref}
\usepackage[top=1.8cm,bottom=1.8cm,left=2.2cm,right=2.2cm]{geometry}

\hypersetup{colorlinks=true,linkcolor=black,urlcolor=blue!70!black}

\setlength{\parindent}{0pt}
\setlength{\parskip}{0.45em}
\setlist[itemize]{leftmargin=1.3em,topsep=0.2em,itemsep=0.15em}
\setlist[enumerate]{leftmargin=1.5em,topsep=0.2em,itemsep=0.15em}

% kompaktere Überschriften
\usepackage{titlesec}
\titlespacing*{\section}{0pt}{0.7em}{0.25em}
\titleformat{\section}{\normalfont\large\bfseries}{\thesection.}{0.5em}{}

\newcommand{\code}[1]{\texttt{#1}}

\begin{document}

\begin{center}
  {\Large\bfseries Mini-Exposé: Kontexttransparenz für Coding Agents}\\[0.3em]
  {\small Inspizierbares, graphbasiertes Kontext-Engineering und seine Wirkung auf
  Effizienz und Entwicklerkontrolle \\
  Ausschreibung: \enquote{Mehr Effizienz mit Coding Agents? Kontexttransparenz und
  Kontext-Engineering}}
\end{center}

\vspace{0.2em}
\hrule
\vspace{0.4em}

% ---------------------------------------------------------------------------
\section{Warum: das Problem}

Bei Coding Agents (Claude Code, Cursor) entscheidet die \textbf{Kontextbereitstellung} über
Erfolg oder Halluzination (Gloaguen et al.\ 2026). Zwei Probleme bleiben offen:

\begin{itemize}
  \item \textbf{Ineffizient.} Agents explorieren Codebases reaktiv und rufen systematisch mehr
    Kontext ab als sie nutzen (ContextBench, 2026). Statische Kontext-Dateien (\code{CLAUDE.md})
    helfen kaum: handgeschrieben nur {\raise.17ex\hbox{$\scriptstyle\sim$}}4\,\% bessere
    Lösungsrate, generiert sogar schlechter, beide ${>}\,20\,\%$ teurer (Gloaguen et al.\ 2026).
  \item \textbf{Intransparent.} Löst der Agent falsch, sieht der Entwickler nicht, \emph{welcher
    Kontext} zu \emph{welcher Entscheidung} führte; Kontext ist eine Black Box. Werkzeuge
    setzen erst \emph{während} der Ausführung an (AgentStepper). Ob man die Kontextauswahl
    \textbf{vor} dem Start bewerten und korrigieren kann und ob das die Qualität verbessert,
    ist bislang kaum systematisch untersucht, obwohl Entwickler diese präventive Kontrolle
    nachweislich suchen (Dhanorkar et al.\ 2026).
\end{itemize}

% ---------------------------------------------------------------------------
\section{Was wird Technisch umgesetzt:}

Ein MCP-Server, der den Kontext für eine Aufgabe nicht nur automatisch assembliert,
 sondern auch die Auswahl vor der Agentausführung inspizier- und korrigierbar macht. 
 Eine Codebase wird als Wissensgraph (Neo4j, Tree-Sitter für Python/TypeScript, 
 lokale Embeddings) repräsentiert. Es gibt drei Funktionen.

\begin{enumerate}
  \item \textbf{Context Assembly (hybrid).} Seed-Knoten über semantische Embedding-Ähnlichkeit
    \emph{und} exakte Namens-/Pfad-Treffer; Expansion per Graph-Traversal (strukturell,
    historisch über Git/Issues, testbasiert) zu einem aufgabenspezifischen Kontext-Block.
  \item \textbf{Entscheidungspfad.} Pro Knoten eine nachvollziehbare Begründung: Score und
    Query-Token-Überlappung für Seeds, konkrete Graph-Kante und Traversal-Kette für den Rest.
  \item \textbf{Inspektions- und Korrektur-Interface.} Ein browserbasiertes Review-UI zeigt den
    geplanten Teilgraphen samt Retrieval-Trace. Der Entwickler sieht \emph{warum} jeder Knoten
    drin ist, kann Knoten entfernen/ergänzen, die Anfrage editieren und neu abfeuern: Kontext
    als editierbarer Vorschlag statt Black Box.
\end{enumerate}

Der eigentliche Beitrag ist nicht der Graph (Stand der Technik), 
sondern die Retrieval-Transparenz als Eingriffspunkt sowie deren empirische Messung im Feld. 
Reines Embedding-Retrieval liefert lediglich eine Rangliste. 
Die explizite Graphstruktur begründet die Auswahl und macht sie korrigierbar.

% ---------------------------------------------------------------------------
\section{Forschungsfragen}

\begin{itemize}
  \item \textbf{HF1 (Effizienz):} Verbessert graph-assemblierter Kontext die Effizienz
    (Iterationen, Dauer, Erstkorrektheit) gegenüber reaktiver Exploration, statischer
    \code{CLAUDE.md} und manuell ausgewählten Snippets?
  \item \textbf{HF2 (Transparenz):} Hilft ein expliziter Entscheidungspfad Entwicklern,
    fehlerhafte Kontextauswahlen \emph{vor} dem Lauf zu erkennen?
  \item \textbf{HF3 (Intervention, explorativ):} Führt die Korrektur der Kontextauswahl vor der
    Ausführung zu besseren Agent-Ergebnissen?
\end{itemize}

% ---------------------------------------------------------------------------
\section{Wie: Methodik}

Design Science Research mit iterativen Build-Evaluate-Zyklen, zwei Studien:
\textbf{Studie 1 (HF1)} vergleicht vier Kontextbedingungen (B1 nur Aufgabe, B2 \code{CLAUDE.md},
B3 manuelle Snippets, B4 Graph) über Bugfix-, Refactoring- und Feature-Aufgaben; Metriken werden
automatisch geloggt. \textbf{Studie 2 (HF2/HF3)} ist eine kontrollierte Studie am Review-Interface
mit gezielt verfälschten Kontextauswahlen (seeded errors). \textbf{Feldumgebung:} aktiv
entwickelte Projekte (Python/TypeScript), Pilot ist ein realer Webshop; der Wert von
Transparenz ist auf eingefrorenen Benchmark-Snapshots prinzipiell nicht messbar.

\textbf{Umsetzung und Stand (Juni 2026):} Die vollständige Umsetzung des Artefakts ist
konstruktiver Bestandteil der Arbeit. Ein lauffähiger Prototyp (Parser, Graph, hybride
Assembly-Pipeline, MCP-Server \emph{und} ein erstes interaktives Review-Interface) belegt
bereits die Machbarkeit; eine erste Eval zeigt brauchbare Kontextabdeckung. Im Verlauf der
Arbeit werden Artefakt und Interface zur Studienreife ausgebaut und anschließend empirisch
untersucht.

\end{document}
